% Appendix I --- Verification: the inventory behind section 1.1 (pass 6, the author's review of
% section 1, 2026-10-26; records/affine/REVIEW-author-paper-affine.md, comment 2), on the isotropic
% paper's model (paper-iso/appendix-verification.tex; article-kit WRITING.md section 4 and
% PUBLICATION-TEMPLATE.md B.4). Section 1.1 states the trust base: what is machine-checked, what
% is assumed, what is computed. This appendix adds detail to that statement and holds no caveat,
% gap or assumption that 1.1 does not state.
%
% Moved here from 01-introduction.tex by records/affine/edits_p6_verification_appendix.py, which
% copies each paragraph byte for byte and checks that it is identical to the text removed from
% section 1.1: "What is proved", the statements that rest on each cited fact and those that rest
% on none, "Standard theorems", "Review" and "Where to look" with the table of the Lean names and
% the stamped #print axioms block (its stamp and content untouched). New: the opening sentence,
% the heading of the second paragraph and its first sentence (which had closed "What it rests on"
% in 1.1), and the first sentence of "Review", which 1.1 keeps, is not repeated.
%
% The region between the two "trust base" comment lines is read by `linkage check` (LINKAGE.md
% rule 13), as the region of section 1.1 is: statuses, the cited facts the region names (by their
% Lean names in the table and the block) against those its \ref'd statements read, and the stamped
% block.

\section{Verification}\label{app:verification}
% trust base: begin
% Q-0377 (2026-10-08; article-kit PROCESS.md § 7, WRITING.md § 4): rewritten when the
% formalization closed (125 of 125 blueprint nodes \leanok on main at dfafe99), on Paper V's model
% (spatial-hemigroup-scale-space, paper/01-introduction.tex, "What is machine-checked"). It replaces
% "No statement of this version is machine-checked" and the four kinds of statement (proved in full,
% sketched, restated, cited). The statuses are the blueprint's. The cited facts each statement rests
% on are what Formalization/CIAxiomGuard.lean prints at dfafe99 (exit code 0, no sorryAx), mapped to
% the ledger through blueprint/trust-boundary.txt: C1 is A6, C2 is A3, C3 is A5, C4 is A9. The lists
% by section are the previous version's, checked against that output; the one addition is the
% isotropic classification, restated in section 2.3, whose formal proof reads C2 (ii)-(iii) at d = 1
% since Q-0365. No statement of this paper read A1, A2, A4, A8 or A10 then; round 3 added C5, which is
% A2, for prop:wedge-axis-profiles(2). Everything between the two
% "trust base" comment lines is read by `linkage check` (LINKAGE.md rule 13): statuses, the ledger
% entries named against those the \ref'd statements read, and the stamped #print axioms block
% (`--fresh-axioms` reruns it). The axioms of the line are not written as (A1)-(A8) inside the
% region, where the checker would read them as ledger entries.
This appendix gives the detail behind \S\ref{sec:trust-base}: what is proved, which statements rest
on which cited fact, the standard theorems the proofs take from Mathlib, the review of the formal
statements, and where the formal development is.

\paragraph{What is proved.} The theorems are machine-checked: the single step and the rigidity, on
kernels and on operators (Theorems~\ref{thm:rank-deficient-step}, \ref{thm:affine-rigidity}
and~\ref{thm:affine-rigidity-operators}), the infinitely divisible increments
(Theorem~\ref{thm:cascade-infinitely-divisible}), the orders of cone fields
(Theorem~\ref{thm:cone-orders}), and the theorem on the general linear group along the chain of
cones (Theorem~\ref{thm:wedge-axis}), whose complete proof Appendix~\ref{app:wedge-axis} collects. So is
every other definition, lemma and proposition of \S\S\ref{sec:classical}--\ref{sec:landscape} and of
the appendices. A proof
marked as a sketch shows the decisive steps and leaves routine verifications out. Twenty statements
have a sketch only: Propositions~\ref{prop:affine-orbits},
\ref{prop:rational-symbols}, \ref{prop:similarity-stable-order-kernels},
\ref{prop:isotropic-corner}, \ref{prop:glplus-families}, \ref{prop:tempered-step-cone}(2)--(4),
\ref{prop:matern-negative-mass}, \ref{prop:fixed-frame-families}, \ref{prop:similarity-ray-families},
\ref{prop:diagonal-ray-families}, \ref{prop:diagonal-stable-order-families},
\ref{prop:similarity-stable-family}, \ref{prop:similarity-orbit-family},
\ref{prop:similarity-general-sandwich}, \ref{prop:similarity-general-bounded-shape},
\ref{prop:similarity-degenerate-family}, \ref{prop:similarity-line-family} and
\ref{prop:similarity-general-index-zero}, and Lemmas~\ref{lem:dilation-ratio}
and~\ref{lem:parameter-space}. Their complete proofs are in the LaTeX sources deposited with this
version (see the code availability statement at the end), whose index,
\texttt{records/affine/INDEX-complete-proofs.md}, names, for each statement
number, the file and the place of its complete proof. A sketched statement is machine-checked like
every other.
% Round-1 external review (Gemini, section 5 item 5; PLAN-review-response-paper.md, D2): the sketched
% statements are listed so that the deposit is a citable address; Propositions 4.15 and 4.19 are
% printed in full in Appendix B.

\paragraph{The cited facts, statement by statement.} The statements whose formal proofs rest on the
five cited facts of \S\ref{sec:trust-base} are these, section by section.
\begin{itemize}
\item \S\ref{sec:line}. C2, its clauses (ii) and (iii) at $d = 1$: the isotropic classification,
Theorem~\ref{thm:isotropic-classification}, restated from the isotropic paper.
\item \S\ref{sec:classical}. C1: Theorem~\ref{thm:cascade-infinitely-divisible}. No other statement
of \S\ref{sec:classical} or of Appendix~\ref{app:classical} rests on one.
% CHANGED Q-0281 wave 2, group H: "and through Theorem cone-orders(4) Theorem wedge-axis" dropped
% from C4 (the entry-free route, the author's decision on W2F-1).
\item \S\ref{sec:landscape}. C2: in \S\ref{sec:affine-stable-cone},
% Round 3 (content review), B.S7: only the increments of 4.7 spend C2.
Lemma~\ref{lem:elliptical-stable}, Proposition~\ref{prop:stable-step-cone}, for the increments of its
clauses (1) and (2), and Propositions~\ref{prop:stable-exponent-cone}, \ref{prop:anisotropy-budget}, \ref{prop:orbit-step-cone},
\ref{prop:bounded-profile-rays} and~\ref{prop:similarity-stable-order-families}; in
\S\ref{sec:smaller-groups}, Proposition~\ref{prop:similarity-stable-order-kernels}; and
Theorem~\ref{thm:wedge-axis} and Proposition~\ref{prop:wedge-axis-profiles}. C3:
Proposition~\ref{prop:anisotropy-budget}, for its clause (1), and through it
Theorem~\ref{thm:wedge-axis}. C1: Proposition~\ref{prop:wedge-axis-profiles}, and
Proposition~\ref{prop:similarity-stable-order-kernels}, for its clause (1). C5:
Proposition~\ref{prop:wedge-axis-profiles}(2). C4: in \S\ref{sec:invariant-wedges}, Lemma~\ref{lem:spectral-cones}, for its
clause (2) and the products in its clause (5); and through it Theorem~\ref{thm:cone-orders}, for its
clauses (1), (2) and (4).
\item The supplementary appendices. C2: Proposition~\ref{prop:tempered-step-cone}
(Appendix~\ref{app:orbits}); Proposition~\ref{prop:diagonal-stable-order-families}
(Appendix~\ref{app:smaller-groups}); Proposition~\ref{prop:one-parameter-families}
(Appendix~\ref{app:linear-flow}); and the eight propositions of Appendix~\ref{app:similarity},
\ref{prop:similarity-stable-family}, \ref{prop:similarity-orbit-family},
\ref{prop:similarity-general-sandwich}, \ref{prop:similarity-general-bounded-shape},
\ref{prop:similarity-degenerate-family}, \ref{prop:similarity-line-family},
\ref{prop:similarity-general-index-zero} and~\ref{prop:similarity-cascade-families}. C1:
Proposition~\ref{prop:similarity-cascade-families}.
\end{itemize}
These are machine-checked on Lean's logical axioms alone, with no cited fact:
\begin{itemize}
\item the isotropic families read on the multiples of the identity
(Proposition~\ref{prop:isotropic-corner});
\item the passage between operators and kernels (Lemmas~\ref{lem:operators-to-kernels}
and~\ref{lem:kernels-to-operators});
\item the rigidity on kernels and on operators (Theorems~\ref{thm:affine-rigidity}
and~\ref{thm:affine-rigidity-operators}), with the single step and the lemmas behind it, and where
continuity enters it (Proposition~\ref{prop:affine-rigidity-boundary});
\item the dictionary between invariant orders and step sets (Proposition~\ref{prop:invariant-orders}),
and the stable orders read in the Loewner order (Proposition~\ref{prop:stable-order-generated}).
\end{itemize}

% CHANGED Q-0281 wave 2: blind review W2C-1 (the author's decision of 2026-10-05: Mathlib's route).
% The proof of thm:cascade-infinitely-divisible(1) now uses Levy's continuity theorem on R^d and the
% portmanteau theorem; they are named here as standard theorems, apart from C1-C4 and the ledger.
% The author's decision of 2026-11-03 (notes/HANDOFF.md, "Waiting on the author"): Prokhorov's theorem,
% used by the proof of lem:operators-to-kernels(1) and taken from Mathlib, is named here too.
\paragraph{Standard theorems.} The proofs use the standard theorems of measure theory, Fourier
analysis, weak convergence and elementary analysis (Fekete's lemma on superadditive functions
among them) without citing them, as they use Fubini's theorem. Three of them are
named here because they are limit theorems of probability, like the cited facts: Prokhorov's
theorem, \Levy's continuity theorem on $\RR^d$ and the portmanteau theorem. The proof of
Lemma~\ref{lem:operators-to-kernels} uses the first, to obtain the kernel of an operator as a weak
limit. The proof of Theorem~\ref{thm:cascade-infinitely-divisible} uses the other two, to show that
the small steps of a cascade form a null array. All three are proved in Mathlib, and the formal
proofs take them from it, whereas they take each cited fact as an assumption.

% Q-0382 (2026-10-11): the verdict of the fidelity review (PROCESS.md § 6, Q-0376;
% blueprint/REVIEW-fidelity.md, "Verdict" and its findings ledger: 0 claim-changing rows, one
% statement-tightening row, R1, a \lean tag that no paper prints, 18 note-only), replacing the
% sentence that said the review was in progress. The witnesses are its section "Witnesses (P1)",
% each on Lean core. The one rewording is Q-0381's (the review's paper-facing point 6), already in
% the shared statement.
\paragraph{Review.} Each formal statement
was read against the statement printed here, in search of one that is weaker than printed, holds
vacuously, or is true only by a convention of the proof assistant, and the review found none. It
found that each cited fact states what its page supports, with the departures recorded in
\S\ref{sec:cited-facts}. For each theorem named under ``What is proved'', a formal example that rests
on no cited fact shows that the hypotheses can be met.

\paragraph{Where to look.} The formal development, its ledger of the cited facts with the
page-verified citation of each, and the LaTeX sources of the statements and their proofs, in which the
declaration behind each numbered statement is named in a \texttt{\char`\\lean} tag on that statement
and the complete proof of every sketched statement is given, are in the repository accompanying this
paper, \url{https://github.com/danielfagerstrom/spatial-hemigroup-scale-space-affine}, at its tag \texttt{v0.1}, as the code availability statement at
the end says.
The toolchain is pinned in-tree: Lean~4 \texttt{v4.31.0}, Mathlib \texttt{v4.31.0}, the shared
library \texttt{ScaleSpaceCore} at its tag \texttt{v0.5.0}. The build configuration requires these
two packages and no other. \texttt{lake build} in \texttt{Formalization/} reproduces the verification, and
\texttt{lake env lean CIAxiomGuard.lean} prints the axioms of every named declaration, to be read
against the file \texttt{blueprint/trust-boundary.txt} beside it. In that file the five cited facts
are these axioms, each in the namespace \texttt{AffineHemigroup}.
\begin{center}
\small
\begin{tabular}{@{}ll@{\qquad}ll@{}}
C1 & \texttt{khintchine\_null\_array} & C2(iii) & \texttt{levy\_khintchine\_converse} \\
C2(i) & \texttt{levy\_khintchine\_exists} & C3 & \texttt{operatorSelfDecomposable\_infinitelyDivisible} \\
C2(ii) & \texttt{levy\_khintchine\_unique} & C4 & \texttt{gelfand\_naimark\_logMajorized} \\
& & C5 & \texttt{sato\_polar\_criterion}
\end{tabular}
\end{center}
On the declarations of Theorems~\ref{thm:affine-rigidity}, \ref{thm:affine-rigidity-operators},
\ref{thm:cascade-infinitely-divisible}, \ref{thm:cone-orders} and~\ref{thm:wedge-axis} and of
Proposition~\ref{prop:wedge-axis-profiles}, it prints Lean's three logical axioms
(propositional extensionality, \texttt{propext}; the axiom of choice, \texttt{Classical.choice}; and
the soundness of quotient types, \texttt{Quot.sound}) and names of the cited facts, and
nothing else:
% printed at fdf0c61 (2026-11-04)
% Q-0382: re-stamped after the fidelity review (Q-0376) and Q-0381, Q-0407 landed. After a full
% `lake build` at 91f39d8 (the Lean sources of main there), `linkage check --fresh-axioms` reran
% `#print axioms` on these declarations and found the block equal to Lean's output, so no line
% of it changes. First printed at dfafe99 (2026-10-08) by `lake env lean CIAxiomGuard.lean`.
% wedge_axis_cell added when the Lean of prop:wedge-axis-profiles landed (2d7a4a6),
% `linkage check --fresh-axioms` comparing the whole block with Lean's output again; compared again at
% fdf0c61, after the author's decision D16 dropped two hypotheses of wedge_axis_cell, and no line changed.
\begin{verbatim}
'AffineHemigroup.affine_rigidity' depends on axioms: [propext,
 Classical.choice,
 Quot.sound]
'AffineHemigroup.affine_rigidity_operators' depends on axioms: [propext,
 Classical.choice,
 Quot.sound]
'AffineHemigroup.gw_family_infinitelyDivisible' depends on axioms: [propext,
 AffineHemigroup.khintchine_null_array,
 Classical.choice,
 Quot.sound]
'AffineHemigroup.causalLE_order' depends on axioms: [propext,
 AffineHemigroup.gelfand_naimark_logMajorized,
 Classical.choice,
 Quot.sound]
'AffineHemigroup.wedge_axis_row' depends on axioms: [propext,
 AffineHemigroup.levy_khintchine_converse,
 AffineHemigroup.levy_khintchine_exists,
 AffineHemigroup.levy_khintchine_unique,
 AffineHemigroup.operatorSelfDecomposable_infinitelyDivisible,
 Classical.choice,
 Quot.sound]
'AffineHemigroup.wedge_axis_cell' depends on axioms: [propext,
 AffineHemigroup.khintchine_null_array,
 AffineHemigroup.levy_khintchine_converse,
 AffineHemigroup.levy_khintchine_exists,
 AffineHemigroup.levy_khintchine_unique,
 AffineHemigroup.sato_polar_criterion,
 Classical.choice,
 Quot.sound]
\end{verbatim}
% trust base: end

